\documentclass{report}
\usepackage{amsmath,amssymb}
\setlength{\parindent}{0mm}
\setlength{\parskip}{1em}
\begin{document}
\begin{center}
\rule{6in}{1pt} \
{\large Daniel F Martin \\
{\bf Solver Challenges in AMR Ice Sheet Modeling }}

Lawrence Berkeley National Laboratory \\ Mail Stop 50A-1148 \\ 1 Cyclotron Rd \\ Berkeley \\ CA 94720
\\
{\tt dfmartin@lbl.gov}\\
Mark Adams\\
Stephen Cornford\end{center}

Ice sheets are a prime candidate for adaptive mesh refinement (AMR) due
to the combination of localized regions which demand very fine spatial
resolution, like ice streams and grounding lines, combined with large
quiescent regions where such fine resolution represents a waste of
computational effort.
BISICLES is a scalable ice-sheet model built on the Chombo framework
which uses block-structured AMR to dynamically refine the computational
mesh where needed to resolve the wide range of dynamical scales found in
ice sheets.
We use a formulation of the momentum balance based on the
vertically-integrated treatment of Schoof and Hindmarsh, which produces a
nonlinear coupled elliptic system which must be solved repeatedly to
compute the ice velocity field.
The combination of local refinement and the particular non-Newtonian
character of ice produces unique challenges solving this system of
equations. We have adopted a hybrid Picard-JFNK approach to the nonlinear
system.
While geometric multigrid (GMG) is a natural fit with AMR methods, it has
struggled with the linear systems some of our applications produce,
particularly in the case of marine ice sheets, leading us to explore
applications of algebraic multigrid (AMG) schemes to AMR.


\end{document}
